% Brief insertion fragment. Development and provenance in the same folder.
\subsubsection{Arithmetic memory of the multiplicative channel}
\label{prd-arit:seccion}

Preservation of the regional signature extends to the full
arithmetic evaluation of APP. For an oriented position
\(x=(i,j,d)\in\{1,\ldots,9\}^2\times\{N,E,S,O\}\), let
\(\chi(x)=(a,b,\epsilon)\), where \(a\) is the fixed coordinate,
\(b\) the moving coordinate, and \(\epsilon\) its orientation:
\[
 \chi(i,j,N)=(j,i,-1),\quad \chi(i,j,S)=(j,i,+1),\qquad
 \chi(i,j,E)=(i,j,+1),\quad \chi(i,j,O)=(i,j,-1).
\]
The axis label \(\iota\in\{i,j\}\) preserves which of the
two original coordinates moves. Writing
\(\langle b\rangle_9=1+((b-1)\bmod9)\), the operators and
evaluation are
\begin{align}
 \overline P(a,b,\epsilon)&=(a,\langle b+\epsilon\rangle_9,\epsilon),
 &\overline H(a,b,\epsilon)&=(a,b,-\epsilon),\label{prd-arit:PH}\\
 n(a,b)&=ab,&r(a,b)&=\langle ab\rangle_9,\qquad
 q(a,b)=\frac{ab-r(a,b)}9.\label{prd-arit:rq}
\end{align}

\begin{proposition}[Exact representation of arithmetic transport]
\label{prd-arit:minimal}
The identities \(\chi P=\overline P\chi\) and
\(\chi H=\overline H\chi\) hold. The representation \(\chi\)
has \(162\) states and is the coarsest of the representations
stable under \(P,H\) that preserve the integer product or,
equivalently, the pair \((r,q)\). Additionally preserving \(\iota\)
recovers the \(324\) original oriented positions.
\end{proposition}
\begin{proof}
Advance modifies only \(b\), and the half-turn only
\(\epsilon\), proving the two transport identities.
The nine observations \(n_k=n(\overline P^k\chi(x))\),
\(0\leq k<9\), run through \(a,2a,\ldots,9a\). Thus one recovers
\(a=\min_k n_k\), \(b=n_0/a\), and the sign through
\((n_1/a-n_0/a)\bmod9\in\{1,8\}\). Every stable equivalence
that preserves \(n\) preserves these nine observations and therefore
distinguishes all states \((a,b,\epsilon)\).
The identity \(n=r+9q\) proves equivalence of the two
evaluations. Each fiber of \(\chi\) contains the two transposed
routes: \((i,j,d)\mapsto(j,i,d^\top)\), with
\(N^\top=O,O^\top=N,E^\top=S,S^\top=E\). The axis label distinguishes
these two preimages and allows position and direction to be reconstructed.
\end{proof}

The TRIT selector activates the product at \(t\equiv2,5,8\pmod9\),
before each advance. The half-turn acts at the end of every
\(27\) instants. If \(y_{t-1}\) denotes the preceding arithmetic
state, the increment and its accumulated memory are determined by
\begin{equation}
 \eta_q(t)=\mathbf1_{\{t\equiv2\ (3)\}}q(y_{t-1}),\qquad
 \lambda_q(t)=\lambda_q(t-1)+\eta_q(t),\quad \lambda_q(0)=0.
 \label{prd-arit:incremento}
\end{equation}
The bound \(0\leq q(a,b)\leq a-1\) holds. To invert an
update, one first undoes the corresponding half-turn,
then the advance, and finally subtracts the quotient recovered from the
preceding state. The record retains phase, axis, event order, and
memory; \(\lambda_q\) is an additive cocycle of this transport.

Each block of \(27\) contains nine advances and traverses the
moving coordinate once. Hence the return of position and orientation
at \(54\), and of phase at \(108\), coexist with
\begin{equation}
 \lambda_q(108N)=NQ_\times^+(a),\qquad
 Q_\times^+(a)=\frac{180a-4\sum_{b=1}^9r(a,b)}9,
 \qquad N\geq0.
 \label{prd-arit:retorno}
\end{equation}
Indeed, each block sums the products \(a,2a,\ldots,9a\);
summing \(n=r+9q\) over the four blocks gives the formula. If
\(a\) is coprime to \(9\), the residues sum to \(45\), and
\(Q_\times^+(a)=20(a-1)\). In window \(j\), counted from
\(j=0\), the relative calendar sign is
\(\sigma_j=(-1)^{\lfloor j/3\rfloor}\), and the effective
orientation is \(\epsilon_j=\epsilon_0\sigma_j\).
The charge signed by \(\sigma_j\) vanishes; the charge weighted
by the effective orientation likewise vanishes after multiplication by
\(\epsilon_0\). The ledger preserves the initial orientation, the accumulated
charge, and its increments. These quantities belong
to the multiplicative cursor; the readers of a trajectory that
alternates between both sheets act on the record of that trajectory.
